\section{Infrastructure}

\begin{qfdefinition}[Infrastructure: features of the assets, revenue models and stages, and the investment characteristics of infrastructure]\label{def:infrastructure-characteristics}
\qffield{Statement} \emph{Features.} Infrastructure assets provide essential services through long-lived physical structures with a high construction cost, a long construction period and lives of several decades: transport (toll roads, ports, airports, railways), regulated utilities and energy networks, and communications (fibre, towers) \citep[p.~97]{larreur-structured-finance-2021}; \citep[pp.~141--142]{demaria-private-markets-2021}. Many enjoy a local or national monopoly and are regulated or contracted, which makes them less sensitive to the economic cycle (Demaria, p.~141). Social infrastructure (hospitals, schools, prisons) is often delivered through public--private partnerships (Larreur, pp.~114--116); Demaria and co-authors exclude it as closer to companies or property.

\emph{Revenue models} (Larreur, pp.~110--116):
\begin{itemize}
\item \emph{long-term purchase contracts} (offtake or power purchase agreements) with a creditworthy buyer;
\item \emph{availability payments} in a PPP, paid by a public authority for keeping the asset available to agreed standards, irrespective of demand, and reduced for poor performance;
\item \emph{regulated returns} set by a regulator on an asset base;
\item \emph{user charges with traffic or merchant risk}, such as tolls, which carry more risk than contracted revenue.
\end{itemize}

\emph{Stages.} Development (permits, equity-funded and mostly by industrial sponsors), construction (\emph{greenfield}) and operation (\emph{brownfield}); partly built or capex-heavy assets are \emph{yellowfield}. Financial sponsors buy mostly once assets are operating and de-risked (pp.~98--99, 121).

\emph{Investment characteristics.} Stable and predictable income, prices often indexed to inflation, and low correlation with other strategies are the claimed benefits (Demaria, p.~142). The risks are regulatory, political and sovereign, construction risk in greenfield projects, high leverage, and the sensitivity of long-dated contracted cash flows to interest rates (Interpretation, from Demaria, p.~142, and Larreur, pp.~149--156). Investments are illiquid and lumpy, and funds have 10--15-year lives (Larreur, p.~122).

\qffield{Practice} \emph{Ashdown Hospital} (original; \hyperref[meth:infrastructure-underwriting]{underwriting}): a 200 PPP hospital earns availability payments of 20 a year, indexed at 2 percent, for 25 years, against operating costs of 6 indexed at 2.5 percent: cash flow available for debt service of 14 in the first year, rising slowly. The investor bears construction, performance and cost risk, not demand risk.

\qffield{Conditions} The features are those of the sources' definitions; classifications of assets such as data centres or social infrastructure differ between investors.

\qffield{Pitfalls} Treating all infrastructure as low risk regardless of its revenue model. Ignoring the interest-rate duration of contracted cash flows.

\qffield{See also} \hyperref[meth:infrastructure-underwriting]{underwriting}; \hyperref[meth:infrastructure-allocation]{allocation}.

\end{qfdefinition}

\begin{qfmethod}[Underwriting infrastructure: investment methods and vehicles, the project life cycle and its debt and equity financing, due diligence on project risks, and valuation]\label{meth:infrastructure-underwriting}
\qffield{Statement} \emph{Methods and vehicles} \citep[pp.~117--131]{larreur-structured-finance-2021}: direct equity stakes in project companies, alone or with industrial sponsors; closed-end funds with 10--15-year lives, which sell assets to one another at the end of their lives; listed infrastructure companies (such as airports) and listed funds; yieldcos holding operating renewable assets; and debt, through bank loans or infrastructure debt funds.

\emph{Project finance and the life cycle} (pp.~97--103, 156--176). An infrastructure project is typically financed through a special purpose vehicle with non-recourse debt: lenders are repaid only from the project's cash flow and take security over its shares and assets. In development, sponsors fund permits with equity. In construction, debt and the rest of the equity are drawn, often as share capital plus shareholder loans. In operation, cash flow pays operating costs, then debt service, then the sponsors. Construction loans may be mini-perms that must be refinanced.

\emph{Debt sizing.} Lenders size debt so that net operating cash flow covers debt service by a target \emph{debt service coverage ratio} (DSCR $=$ cash flow over principal plus interest), with lower lock-up and default levels. The \emph{loan life cover ratio} compares the present value of cash flow over the loan's life, plus reserves, with the debt. A maximum gearing (for example 80/20) caps debt even when cash flows would support more. Higher-risk projects, such as those with merchant risk, need higher DSCRs (pp.~159--160, 169--170).

\emph{Due diligence} covers the risks and who bears them: construction (fixed-price construction contracts), resource, credit of the offtaker or authority, market, interest-rate and currency, operating, technological, political, environmental and force majeure (pp.~149--156). \emph{Valuation} discounts the equity cash flows at a rate reflecting the phase (higher in construction) and the revenue model (Interpretation).

\qffield{Practice} \emph{Ashdown Hospital PPP} (original). Construction cost 200 over two years; availability payments 20 indexed at 2 percent, operating costs 6 indexed at 2.5 percent, for 25 years; senior debt at 5 percent over 22 years, sculpted to a 1.30 DSCR, with gearing capped at 85 percent:
\begin{itemize}
\item the DSCR supports 165.8 of debt, below the 170 cap, so the DSCR binds (gearing 82.9 percent); the loan life cover ratio at financial close is 1.30, as sculpting implies;
\item equity of 34.2, paid during construction, earns 10.70 percent from distributions after debt service and from the three years after the loan;
\item a 10 percent availability deduction for three years cuts the DSCR to 1.11, just above a 1.10 lock-up;
\item with demand risk instead, at a 1.60 DSCR, the same cash flows support only 134.7 of debt and the equity return falls to 8.47 percent.
\end{itemize}
\emph{Decision}: the availability structure suits a core infrastructure investor; a demand-risk bid on the same assumptions must either cut the price or accept a lower return.

\qffield{Conditions} The model ignores tax, reserve accounts and lifecycle costs; Larreur's account is a practitioner's description of market practice.

\qffield{Pitfalls} Underwriting greenfield projects at brownfield returns. Ignoring refinancing risk in mini-perm structures. Assuming availability payments are free of performance deductions.

\qffield{See also} \hyperref[def:infrastructure-characteristics]{characteristics}; \hyperref[meth:private-debt-credit-analysis]{credit analysis}.

\end{qfmethod}

\begin{qfmethod}[Infrastructure in a strategic allocation: risk and return against other investments, its roles, and a recommendation]\label{meth:infrastructure-allocation}
\qffield{Statement} \emph{Return evidence.} Twenty-nine infrastructure funds of the 1998--2008 vintages earned a pooled net IRR of 6.4 percent, a multiple of 1.44 and a PME of 1.40 against the MSCI World; quartiles ranged from 1.6 to 9.7 percent and the bottom 5 percent of funds earned $-15.8$ percent \citep[p.~142]{demaria-private-markets-2021}. Expected net returns are about 6--8 percent for core, 10--12 percent for value-added and above 14 percent for opportunistic infrastructure (p.~142, citing Rosenbaum).

\emph{Roles} (Interpretation, from the sources above): long-dated, often inflation-indexed income that matches long liabilities; diversification from economic cycles where revenue is contracted or regulated; and a return premium for illiquidity and construction or demand risk.

\emph{Against other investments.} Contracted infrastructure behaves partly like long inflation-linked debt and partly like equity in the operator; demand-exposed assets behave more like equity. Its interest-rate sensitivity means it can fall when rates rise, as bonds do. Listed infrastructure companies add business risk beyond the original asset, so investors require higher returns from them \citep[p.~122]{larreur-structured-finance-2021}; they are more liquid but trade with equity markets (Interpretation).

\qffield{Practice} \emph{An insurer with long annuity liabilities} (original): allocate 5 percent to core brownfield infrastructure equity and 5 percent to infrastructure debt. Availability and regulated assets, such as Ashdown's 10.7 percent equity return on an 83 percent geared project, provide the inflation-linked cash flows. Greenfield and merchant risk are limited to a small opportunistic sleeve, because their equity returns fall sharply when lenders demand higher cover. Commitments are spread over vintages, and the infrastructure return is compared with the PME, not the absolute IRR.

\qffield{Conditions} The return data are Cambridge Associates statistics for a small sample of early vintages, quoted by Demaria and co-authors.

\qffield{Pitfalls} Treating infrastructure as a bond substitute without its equity and regulatory risk. Buying listed infrastructure for diversification from equities.

\qffield{See also} \hyperref[meth:infrastructure-underwriting]{underwriting}; \hyperref[meth:real-assets-allocation]{real assets}; \hyperref[meth:alternatives-allocation-approaches]{allocating to alternatives}.

\end{qfmethod}
